% Created 2026-09-29 Tue 15:49
% Intended LaTeX compiler: xelatex
\documentclass[11pt]{article}
\usepackage{graphicx}
\usepackage{longtable}
\usepackage{wrapfig}
\usepackage{rotating}
\usepackage[normalem]{ulem}
\usepackage{capt-of}
\usepackage{hyperref}
\usepackage[margin=1in]{geometry}
\author{Christopher Wiggins}
\date{\today}
\title{Proposal Experiment Workflow}
\hypersetup{
 pdfauthor={Christopher Wiggins},
 pdftitle={Proposal Experiment Workflow},
 pdfkeywords={},
 pdfsubject={},
 pdfcreator={Emacs 30.2 (Org mode 9.7.11)}, 
 pdflang={English}}
\begin{document}

\maketitle
\clearpage
\section{Experimental Workflow and Testbed Methodology}
\label{sec:orgeada986}

The proposed research will use a controlled, reproducible virtualized
testbed to experimentally evaluate relationships between Linux security
configuration recommendations and defined cybersecurity attack outcomes.
The initial pilot will focus on a single Linux distribution, a single
CIS Benchmark version and profile, one defined attack scenario, and a
small set of selected CIS recommendations. The purpose of the pilot is
not to assume that CIS recommendations prevent particular attacks, but
to determine experimentally whether individual recommendations produce
measurable changes in attack outcomes, operational behavior, and
observable system evidence.

The initial pilot will focus on Ubuntu 24.04 LTS, one identified CIS
Benchmark version and profile, one defined attack scenario, and a
limited number of selected CIS recommendations. The methodology is
designed to determine not only whether CIS hardening as a whole changes
an attack outcome, but also which individual recommendations materially
contribute to that outcome, which depend on other controls, and what
operational consequences accompany those changes.
\section{Virtual Testbed}
\label{sec:org4202cea}
The experimental environment will be hosted on a dedicated Proxmox
Virtual Environment (Proxmox VE) server. Proxmox will provide
virtual-machine isolation, repeatable system restoration, virtual
networking, and snapshot management.

The initial testbed will contain two virtual machines connected through
an isolated virtual network:

\begin{verbatim}
                         PROXMOX HOST
                              |
                 isolated Linux bridge vmbr1
                    NO physical interface
                              |
              +---------------+---------------+
              |                               |
       +------v---------+              +------v---------+
       | ATTACKER VM    |              | TARGET VM      |
       | Debian/Kali    |              | Ubuntu 24.04   |
       |                |              |                |
       | 10.10.10.10/27 |              | 10.10.10.20/27 |
       +----------------+              +------+---------+
                                              |
                                       CIS configuration
                                       OpenSCAP
                                       auditd
                                       journald
                                       ssh logs
                                       ansible automation
\end{verbatim}

\textbf{1. Target VM:} An Ubuntu 24.04 LTS Server system representing the
system being evaluated. The target will contain OpenSSH, OpenSCAP, Linux
Audit, system logging facilities, and the configuration and security
assessment tools required by the experiment.

\textbf{2. Attacker/Test VM:} A separate Kali Linux virtual machine used to
conduct authorized attack simulations, legitimate-access tests, and
external observations of the target.

The experimental network will have no direct connection to production
university systems or public networks during attack testing. All
accounts, credentials, addresses, and services used during
experimentation will be created specifically for the research
environment.

A third data-collection or orchestration system may be introduced as the
project develops. However, the initial pilot will rely primarily on
standard Linux logging mechanisms and experiment scripts so that the
feasibility of the study does not depend on the availability of a
specialized provenance platform.
\section{Baseline Target Configuration}
\label{sec:orgba9cf8b}
Two principal reference states will be maintained.

The first will be a stock Ubuntu baseline, representing the operating
system in a minimally modified or default configuration appropriate to
the selected attack scenario. This system will provide an external
comparison for determining the overall effect of CIS hardening.

The second and more important experimental reference will be a CIS Level
2 hardened golden image. This system will be created from the same
Ubuntu installation, hardened according to the selected CIS Benchmark
version and profile, and evaluated to verify that the intended hardened
configuration has been successfully applied.

The resulting relationship will be:


\textbf{Stock Ubuntu → Apply CIS Level 2 Hardening → Verify Configuration → CIS
Level 2 Golden Image}


Once the hardened system has been validated, a Proxmox snapshot will
preserve it as the canonical hardened reference state. This snapshot
will not be directly modified during experimentation. Instead,
experimental conditions will be derived from the validated golden image.

This design ensures that experimental systems share the same operating
system version, kernel, packages, users, services, instrumentation, and
unrelated security settings. Experimental differences can therefore be
limited to the selected CIS recommendations under investigation.
\section{Control Ablation Method}
\label{sec:orgcad0164}
The primary experimental method will use systematic control ablation. In
this approach, the experiment begins with the complete validated
hardened configuration and selectively removes or weakens individual CIS
recommendations.

For example, if two recommendations are designated Control A and Control
B, the primary experimental conditions will be:


\begin{center}
\begin{tabular}{llll}
Condition & Control A & Control B & Remaining Controls\\
\hline
H0 - hardened & enabled & enabled & enabled\\
H1 - A ablated & less restrictive & enabled & enabled\\
H2 - B ablated & Enabled & less restrictive & enabled\\
H3 - A+B ablated & Less restrictive & less restrictive & enabled\\
\end{tabular}
\end{center}


A separate stock Ubuntu condition will provide an external comparison:


\begin{center}
\begin{tabular}{ll}
Condition & Description\\
\hline
U0 - stock baseline & default ubuntu configuration\\
H0 - hardened ref. & validated CIS level 2 configuration\\
H1 & CIS Level 2 except control A\\
H2 & CIS level 2 except Control B\\
H3 & CIS level 2 except Controls A and B\\
\end{tabular}
\end{center}


This design provides two complementary forms of evidence.

First, comparing U0 with H0 measures the aggregate effect of the
hardened configuration against the selected attack.

Second, comparing H0 with H1, H2, and H3 measures the marginal and
conditional contribution of individual CIS recommendations within an
otherwise hardened environment.

This distinction is important because security controls may overlap.
Removing one control may produce no observable change because another
hardened configuration still prevents the attack. Such a result will not
automatically be interpreted as evidence that the removed control is
ineffective. Instead, it may indicate that the control's effect is
redundant or conditional under the tested configuration.
\section{CIS Recommendation Selection and Mapping}
\label{sec:org4ca9d52}
Candidate recommendations will be selected from a specific CIS Benchmark
version and profile applicable to the target operating system. Each
candidate recommendation will be mapped to a proposed attack
relationship before experimentation.

For each recommendation, the research will document the exact CIS
recommendation identifier and title, CIS rationale, affected
configuration parameter, recommended secure state, deliberately less
restrictive comparison state, hypothesized relationship to the attack,
expected attack outcome, expected operational effect, supporting
technical evidence, and confidence in the proposed relationship.

MITRE ATT\&CK will be used to characterize the attack behavior
represented by the experiment. CWE, CVE, and CVSS information will be
incorporated only when a defensible relationship exists. A CIS
recommendation addressing a configuration weakness will not
automatically be associated with a CVE or CVSS score.

This mapping stage produces a testable hypothesis rather than assuming
that the relationship described by the benchmark is experimentally
established.
\section{Experimental Run Procedure}
\label{sec:org71985f1}

Each experiment run will follow the same controlled sequence:

\begin{enumerate}
\item Restore the appropriate validated reference snapshot.
\item Assign a unique experiment identifier.
\item Apply the control-ablation configuration required by the assigned
experimental condition.
\item Verify that the selected controls are in the intended state.
\item Verify that unrelated hardened controls remain unchanged.
\item Restart or reload affected services where necessary.
\item Confirm that required services are operational.
\item Perform the approved legitimate-use test.
\item Execute the authorized attack scenario from the attacker/test VM.
\item Determine the attack outcome using an independent ground-truth
mechanism.
\item Collect relevant authentication, system, audit, and service logs.
\item Record selected operational and performance measurements.
\item Preserve all configuration, attack, ground-truth, and logging
artifacts.
\item Reset the target before beginning the next experimental run.
\end{enumerate}

Experimental conditions will be executed in randomized order when
practical to reduce the possibility that temporal or environmental
factors systematically affect one condition.

Each primary condition will be repeated at least five times. Where
automation and available time permit, ten repetitions per condition will
be conducted.
\section{Configuration Verification}
\label{sec:orgd1138b0}
Configuration state will be treated as an experimental measurement
rather than assumed from successful execution of a configuration script.

OpenSCAP and direct configuration checks will be used to establish that
the system was in the intended state at the time of each experiment.

For the full hardened reference, the experiment will verify that the
selected CIS Level 2 configuration is present.

For an ablation condition, verification will establish both that:

\begin{enumerate}
\item the selected recommendation was successfully changed to the intended
comparison state; and
\item unrelated controls remained in their hardened states.
\end{enumerate}

The experimental chain therefore becomes:


\textbf{CIS Recommendation → Hardened Configuration → Controlled Ablation →
Configuration Verification → Attack Execution → Observed Outcome}


This approach strengthens causal interpretation by minimizing
configuration differences between experimental conditions.
\section{Independent Ground Truth}
\label{sec:org3a435f9}
Attack success or failure will be established independently of the
logging mechanisms being evaluated.

For an SSH-based experiment, successful access may be confirmed through
a combination of SSH client exit status, establishment of an
authenticated session, execution of a harmless predefined command, and
creation of an experiment-specific marker on the target.

The experiment will therefore maintain two distinct forms of evidence:

\textbf{Ground Truth:} Independently established evidence describing what
actually occurred.

\textbf{System Evidence:} Authentication, journal, audit, process, and other
system records describing what the target reported.

Separating ground truth from system evidence allows the project to
evaluate both the security effect of a configuration and the accuracy of
the evidence produced during the event.
\section{Legitimate-Use and Operational-Impact Testing}
\label{sec:orgcb5c02e}
Every experimental condition will also be evaluated for its effect on
legitimate system operation.

A dedicated approved test account will perform the intended legitimate
operation under each condition. Measurements will include authentication
success, shell or command execution success, login completion time,
service availability, authentication errors, unexpected connection
termination, required manual intervention, and difficulty associated
with reversing the configuration.

This provides a second experimental dimension:

\textbf{Security Effectiveness versus Operational Cost}

A recommendation that prevents an attack without interfering with
legitimate operation represents a different outcome from a
recommendation that prevents the same attack only by substantially
degrading expected system functionality.

\begin{itemize}
\item Interaction Between Security Controls
\end{itemize}

The hardened-reference methodology explicitly recognizes that CIS
recommendations may not operate independently.

For example, disabling one authentication-related control may not change
an attack outcome if another Level 2 control independently prevents the
same attack. Conversely, disabling two controls simultaneously may
expose an attack path that remains blocked when either control is
enabled.

The H1, H2, and H3 conditions therefore allow the experiment to
distinguish between individual and combined effects.

If removing Control A alone produces no change, removing Control B alone
produces no change, but removing both produces a repeatable change in
attack outcome, the evidence suggests an interaction or redundancy
between the controls rather than demonstrating that either control is
individually ineffective.

This interaction analysis is an important reason for retaining combined
ablation conditions.
\section{Logging and Provenance Evaluation}
\label{sec:org0d4d8a0}
Each experimental run will collect available evidence from SSH
authentication logging, the system journal, Linux Audit, and relevant
process or service events.

The collected evidence will be compared against independently
established ground truth to determine whether:

\begin{itemize}
\item configuration changes were observable;
\item authentication attempts were recorded;
\item successful and unsuccessful access was represented correctly;
\item resulting shell or process activity was visible;
\item expected events were missing;
\item unrelated events were incorrectly associated with the attack;
\item experimental failures could be distinguished from successful security
controls; and
\item the recorded event sequence was sufficient to explain why an attack
succeeded or failed.
\end{itemize}

Additional provenance instrumentation may be incorporated in later
phases of the research, but the initial experiment will establish a
baseline using standard Linux evidence sources.
\section{Experimental Data Management}
\label{sec:orgf8dba1f}
Every experimental run will receive a unique identifier encoding the
platform, attack scenario, condition, and repetition. For example:

\texttt{UBU24-SSH-C2-R04}

Each run will generate both structured measurements and supporting raw
artifacts.

The structured dataset will include the experiment identifier, platform
information, CIS Benchmark and recommendation identifiers, configuration
state, verification results, attack scenario and ATT\&CK mapping, attack
outcome, shell-access status, number of attempts, time to outcome,
legitimate-access results, operational measurements, logging results,
validity of the experimental run, and relevant notes.

Raw artifacts will include OpenSCAP results, configuration-verification
output, authentication logs, system journal records, Linux Audit
records, attack-tool output, legitimate-access results, and other
supporting evidence.

Failed experimental runs will remain part of the research record and
will be explicitly distinguished from attacks that were successfully
blocked by the tested configuration.
\section{Experiment Automation and Reproducibility}
\label{sec:org7e87aad}
The experimental framework will progressively automate the complete test
lifecycle. Scripts will be developed for control modification,
configuration verification, legitimate-use testing, attack execution,
evidence collection, result storage, and environment restoration.

Proxmox snapshots will provide the known starting states required for
repeatable experimentation. The intended automated lifecycle is:


\textbf{Restore → Apply Ablation → Verify → Legitimate Test → Attack →
Establish Ground Truth → Collect Evidence → Store Results → Reset}


This automation will reduce variation between repetitions and make it
possible to execute larger experimental matrices as the project expands.

The complete testbed definition, scripts, configuration states,
experiment metadata, raw evidence, analysis procedures, and
documentation will be maintained so that another researcher can
reproduce the experiment.
\section{Analysis and Classification}
\label{sec:orga3cfb7f}
Results will be compared across the four experimental conditions.
Primary measurements will include attack success rate, furthest attack
progression, number of attempts, time to outcome, legitimate-access
success, operational effects, and completeness and accuracy of system
evidence.

Each tested CIS recommendation will ultimately be classified as:

\begin{itemize}
\item Supported: The recommendation produced a repeatable change in the
defined attack outcome.
\item Unsupported: The recommendation did not produce the hypothesized
change.
\item Conditional: The recommendation affected the outcome only when another
configuration or environmental condition was present.
\item Inconclusive: The available experimental evidence was insufficient to
support a defensible conclusion.
\end{itemize}

Null and unsupported results will be retained because they provide
information about the actual relationship between configuration
recommendations and attack behavior.
\section{Extension to the Larger Research Project}
\label{sec:orgbe969b8}
The pilot testbed is intended to establish the feasibility of a broader
empirical security-configuration research framework. Once the
experimental process is validated, the same methodology can be extended
to additional CIS recommendations, attack behaviors, system
configurations, and Linux environments.

The resulting research pipeline can be represented as:

\textbf{CIS Recommendation → System Configuration → ATT\&CK Behavior → CWE/CVE
Relationship Where Applicable → Controlled Experiment → Attack Outcome →
Operational Impact → Logging/Provenance Evidence → Empirical
Classification}

The long-term objective is therefore not simply to determine whether a
CIS-hardened system is more secure than a default system. Rather, the
project seeks to determine \textbf{which configuration recommendations
contribute to particular defensive outcomes, under what conditions those
contributions occur, what operational costs accompany them, and what
observable system evidence is produced.}

This framework provides a foundation for developing a structured dataset
connecting security configuration guidance with experimentally observed
cybersecurity outcomes. Such a dataset could support future research and
tools that provide administrators with more contextual information than
configuration-compliance status alone, including the likely security
benefit, operational tradeoffs, attack relationships, and confidence
supported by experimental evidence.
\end{document}
